% Generated draft: requires mathematical and source review.
\documentclass[11pt]{article}
\usepackage[margin=25mm]{geometry}
\usepackage{fontspec}
\setmainfont{Latin Modern Roman}
\usepackage{amsmath,amssymb,mathtools,unicode-math}
\setmathfont{Latin Modern Math}
\usepackage{enumitem,needspace,xurl,hyperref}
\setlength{\parindent}{0pt}
\setlength{\parskip}{5pt}
\setlength{\emergencystretch}{3em}
\newcommand{\N}{\mathbb{N}}
\newcommand{\Z}{\mathbb{Z}}
\newcommand{\Q}{\mathbb{Q}}
\newcommand{\R}{\mathbb{R}}
\newcommand{\C}{\mathbb{C}}
\newcommand{\F}{\mathbb{F}}
\newcommand{\A}{\mathbb{A}}
\newcommand{\PP}{\mathbb{P}}
\newcommand{\E}{\mathbb{E}}
\newcommand{\eps}{\varepsilon}
\newcommand{\dd}{\,\mathrm{d}}
\begin{document}
\section*{UnsolvedMath ID: 30000157}
\textbf{Source problem code:} catalogue entry
\subsection{Definitions and mathematical statement}
The source's Conjecture 2 concerns nonintegrable birational maps \(L\) that are not reversible over \(\R\) and have no other symmetry. For such an \(L\), let
\[
 \mathcal L_P=\{q\leq P:q\text{ is prime and }L\text{ has the intended good reduction at }q\}.
\]
Write \(D_{0,q}(x)\) for the quantity denoted \(D_0(x)\) at the prime \(q\). Interpreting the source's average as the unweighted arithmetic average, set, when \(\mathcal L_P\ne\varnothing\),
\[
 \langle D_0\rangle_P(x)=\frac{1}{|\mathcal L_P|}
 \sum_{q\in\mathcal L_P}D_{0,q}(x).
\]
The conjectured limit is
\[
 \text{for every qualifying }L,\qquad
 \lim_{P\to\infty}\langle D_0\rangle_P(x)=x
 \quad\text{for each }x\text{ in the intended domain of }D_0.
\]
The averaging cutoff \(P\) is distinguished from the individual reduction prime \(q\); the source writes \(p\) for the cutoff. No uniformity in \(x\), rate of convergence, or averaging over maps is asserted.

The supplied record does not define \(D_0\), its normalization, the cycle sampling convention or its domain in \(x\). In particular, a statement for all real \(x\) cannot be inferred merely from the displayed right-hand side \(x\); if \(D_0\) is a normalized distribution function, its intended interval matters. The finite sets on which reduced maps act, the treatment of indeterminacy and exceptional points, and the arithmetic coefficient field needed for reduction are also missing. The original wording permits primes at which \(L\) ``can be reduced,'' whereas the cleaned wording requires good reduction; their exact equivalence needs the intended reduction convention. Infinitely many eligible primes and well-defined statistics along them are prerequisites not established by the record.

In a common birational convention reversibility means existence of an invertible change of variables \(J\) satisfying \(JLJ^{-1}=L^{-1}\), with \(J\) defined over \(\R\) in the cleaned wording. Whether \(J\) must be an involution, and which transformations qualify, is unspecified. Nonintegrability and ``no other symmetry'' likewise require definitions. The latter cannot simply mean no commuting birational transformations at all, since powers of \(L\) commute with \(L\); the intended exclusion of trivial symmetries is not supplied.

The original text motivates prime averaging by very long cycles, of order the square of the reduction prime, as a signature of random permutations. This is retained as source motivation, not an additional proved hypothesis. The 2026-08-21 literature triage says later reversible-map work models symmetric maps but does not establish this averaged limit for every nonintegrable nonreversible map. Its open label remains dated metadata.
\subsection{Short English statement}
For every birational map in the stated nonsymmetric, nonintegrable class, the conjecture predicts that the prime-averaged cycle statistic tends to \(x\). The statistic, arithmetic reduction setup and precise symmetry exclusions are missing. Results for reversible maps are reported as related work, not a proof of this conjecture.
\subsection{Sources}
\begin{itemize}
\item Local UnsolvedMath registry, permanent ID 30000157, source code catalogue entry, ``Cycle Statistics of Nonreversible Birational Maps,'' including the original Conjecture 2, prime-average motivation and cleaned statement.
\item \emph{Finite Fields: Theory and Applications (2005)}, Oberwolfach Reports, pp.\ 2913--2970, DOI 10.4171/owr/2004/54: \url{https://doi.org/10.4171/owr/2004/54}. The title uses 2005, whereas the DOI and proposed-year metadata use 2004.
\item Supplied reference ``Roberts--Vivaldi, 2009'': \url{https://arxiv.org/abs/0905.4135}.
\item Literature triage checked 2026-08-21; record updated 2026-09-16. The verification note reports removal of explanatory spillover and normalization of notation. No external source has been visited or independently verified here.
\end{itemize}
\end{document}
